\documentclass[12pt]{article}

\usepackage{sbc-template}
\usepackage{graphicx,url}
\usepackage[english]{babel}
\usepackage[utf8]{inputenc}
\usepackage{enumitem}
\usepackage{booktabs}
\usepackage{amsmath}
\usepackage{amssymb}
\usepackage{xcolor}
\usepackage{xspace}
\usepackage{hyperref}
\hypersetup{hidelinks}
% allow long URLs in footnotes to break at more boundaries
\def\UrlBreaks{\do\/\do\-\do\.\do\_\do\?\do\&\do\=\do\#\do\:}
\Urlmuskip=0mu plus 3mu
% ragged footnotes: avoid underfull justification on long URL labels
\makeatletter
\long\def\@makefntext#1{\parindent 1em\noindent
  \hb@xt@1.8em{\hss\@makefnmark}\raggedright #1}
\makeatother

% --- tight layout knobs ---
\renewcommand{\topfraction}{0.95}
\renewcommand{\bottomfraction}{0.95}
\renewcommand{\textfraction}{0.05}
\setlist{itemsep=1pt,topsep=2pt,parsep=0pt}
\raggedbottom  % do not stretch short pages: avoids underfull \vbox during \output
% Times has no small-caps-italic shape; substitute upright small caps so \textsc in an
% italic context does not emit an "undefined font shape OT1/ptm/m/scit" warning. Declared
% after the font family is loaded.
\AtBeginDocument{\DeclareFontShape{OT1}{ptm}{m}{scit}{<->ssub*ptm/m/sc}{}}

% --- system name macro (single point for anonymization/renaming) ---
\newcommand{\sysname}{\textsc{RAGtrap}\xspace}

% --- positioning-table marks ---
\newcommand{\cmark}{\ensuremath{\checkmark}}
\newcommand{\xmark}{\ensuremath{\times}}

% ======================================================================
% RESULTS MACROS -- every empirical number in the paper body is defined
% in results/macros.tex, generated by scripts/aggregate_results.py from
% results/{real_results,scaling_results,aux_results}.json. Prose never
% hardcodes a number; a data refresh is a one-file regeneration.
% ======================================================================
% ===== AUTO-GENERATED RESULTS MACROS (scripts/aggregate_results.py) =====
% Every empirical number in the paper body is defined here, from results/*.json.
\newcommand{\TopK}{10}
\newcommand{\Nquestions}{100}
\newcommand{\NsuspectsE}{1000}
\newcommand{\NpoisonE}{500}
\newcommand{\NcleanE}{500}
\newcommand{\RTrecall}{1.00}
\newcommand{\RTprecision}{1.00}
\newcommand{\RTfpr}{0.00}
\newcommand{\RTlatPerSus}{101.7}
\newcommand{\RTworkE}{1000}
\newcommand{\RTrecallDriftThree}{0.70}
\newcommand{\RTrecallDriftThreeLo}{0.66}
\newcommand{\RTrecallDriftThreeHi}{0.74}
\newcommand{\RTfnrDriftThree}{0.30}
\newcommand{\RTrecallDriftFive}{0.51}
\newcommand{\RTrecallDriftFiveLo}{0.46}
\newcommand{\RTrecallDriftFiveHi}{0.55}
\newcommand{\RTfnrDriftFive}{0.49}
\newcommand{\BLrecall}{0.96}
\newcommand{\BLrecallLo}{0.93}
\newcommand{\BLrecallHi}{0.97}
\newcommand{\BLprecision}{0.88}
\newcommand{\BLprecisionLo}{0.85}
\newcommand{\BLprecisionHi}{0.91}
\newcommand{\BLfpr}{0.13}
\newcommand{\BLfnr}{0.04}
\newcommand{\BLcalls}{1000}
\newcommand{\BLjudge}{Qwen2.5-3B-Instruct}
\newcommand{\LatSpeedup}{16274}
\newcommand{\RTlatTotal}{101.7}
\newcommand{\BLlatTotal}{1654}
\newcommand{\BLperSus}{1.7}
\newcommand{\FullPassages}{2{,}681{,}468}
\newcommand{\FullChunks}{4{,}364{,}162}
\newcommand{\EdLatency}{69.2}
\newcommand{\EdThroughput}{14453}
\newcommand{\EdRecord}{575}
\newcommand{\RevChunks}{100}
\newcommand{\FullRevMttr}{46.2}
\newcommand{\FullManMttr}{782}
\newcommand{\FullMttrRatio}{16931}
\newcommand{\SmallPassages}{10{,}000}
\newcommand{\SmallMttrRatio}{44}
\newcommand{\EdSig}{64}
\newcommand{\HmacSig}{32}
\newcommand{\HmacLatency}{33.0}
\newcommand{\EdOverHmac}{1.9}
\newcommand{\EThreeDocs}{230}
\newcommand{\PerDocFP}{0.52}
\newcommand{\PerDocFPLo}{0.50}
\newcommand{\PerDocFPHi}{0.55}
\newcommand{\PerDocPurged}{1441}
\newcommand{\PerDocFalse}{751}
\newcommand{\PerChunkFP}{0.00}
\newcommand{\AsrPoint}{98}
\newcommand{\AsrLo}{93}
\newcommand{\AsrHi}{99}
\newcommand{\AsrN}{100}
\newcommand{\AsrTopK}{5}
\newcommand{\AsrGen}{Qwen2.5-3B-Instruct}
\newcommand{\EZeroChunks}{200}
\newcommand{\EZeroPurged}{20}
\newcommand{\PoisonTexts}{5}
\newcommand{\PoisonASR}{90}
% ===== END AUTO-GENERATED MACROS =====

% ======================================================================

\sloppy
\setlength{\emergencystretch}{3em}

\title{\sysname{}: Per-Chunk Signed Provenance with Constant-Time\\
Traceback and One-Command Source Revocation for RAG Pipelines}

\author{Anonymous Author(s)\inst{1}}

\address{Affiliation Anonymized for Double-Blind Review\\
  \email{anonymous@example.org}}

\begin{document}

\maketitle

\begin{abstract}
Retrieval-augmented generation (RAG), which grounds a language model on
passages fetched from an external corpus, is now the dominant enterprise
pattern, yet that corpus is ingested from untrusted sources with no trust
boundary at admission. Corpus poisoning is cheap and effective: published
attacks reach roughly \textbf{\PoisonASR\%} attack success by injecting only
\textbf{\PoisonTexts} malicious passages per target question. Existing defenses
filter at query time or run a post-hoc traceback driven by an iterative
language-model judge, costly and reactive per incident, and keep no provenance
for recovery. We present \sysname{}, an ingestion gate that writes, per chunk,
an Ed25519-signed provenance record natively into the vector store, turning
poisoning traceback into a single constant-time signature lookup and enabling
one-command revocation of a compromised source. On the real PoisonedRAG attack
over Natural Questions, with suspects surfaced by the released dense-retrieval
feedback, \sysname{} matches the published RAGForensics LLM-judge baseline on
attribution but is \textbf{$\sim$\LatSpeedup$\times$} faster at zero per-incident
model cost, and revokes a compromised source in constant time whose advantage
over a full-corpus purge grows to \textbf{\FullMttrRatio$\times$} on the full
\FullPassages{}-passage corpus. We release \sysname{} as a reproducible artifact.
\end{abstract}

\section{Introduction}

Retrieval-augmented generation (RAG), the technique of answering a query by
first retrieving relevant passages from an external document corpus and then
conditioning a large language model (LLM) on them, has become the dominant way
organizations deploy LLMs over their own data. Enterprise RAG adoption rose to
\textbf{51\%} in 2024 from 31\% a year
earlier,\footnote{\href{https://menlovc.com/2024-the-state-of-generative-ai-in-the-enterprise/}{Menlo
Ventures, ``2024: The State of Generative AI in the Enterprise'', 2024.}} and
around \textbf{70\%} of organizations using generative AI now ground their
models with retrieval systems and vector databases rather than off-the-shelf
models alone, a category that grew 377\% year over
year.\footnote{\href{https://www.databricks.com/blog/state-ai-enterprise-adoption-growth-trends}{Databricks,
``State of AI: Enterprise Adoption \& Growth Trends'', 2024.}} The RAG market,
estimated at \textbf{USD 1.2 billion} in 2024, is projected to grow at a
\textbf{49.1\%} compound annual rate through
2030.\footnote{\href{https://www.grandviewresearch.com/industry-analysis/retrieval-augmented-generation-rag-market-report}{Grand
View Research, ``Retrieval Augmented Generation Market Size Report, 2030''.}}
The pattern is consequential precisely because it carries high-stakes content:
government agencies adopt RAG to answer citizens from official documents while
keeping regulated data in place and citing the source of each
answer,\footnote{\href{https://fedtechmagazine.com/article/2025/05/how-rag-boosts-agency-accuracy-security-perfcon}{A.~Slagg,
``How RAG Boosts Agency Accuracy and Security'', FedTech, 2025.}}$^{,}$\footnote{\href{https://statetechmagazine.com/article/2025/02/what-is-rag-perfcon}{StateTech,
``What Is Retrieval Augmented Generation, and How Are State and Local Agencies
Using It?'', 2025.}} so the integrity and traceability of the corpus is a
first-class requirement, not an afterthought.

That corpus is the attack surface. Because the passages are ingested from the
open web and shared document stores with no trust boundary at admission, an
adversary who controls even a few sources can poison the answers. PoisonedRAG
shows that injecting only \PoisonTexts{} malicious passages per target question
into a knowledge base of millions of texts reaches about \PoisonASR\% attack
success, steering the model to an attacker-chosen answer~\cite{poisonedrag}.
Carlini et al. show this is practical at web scale: an attacker could poison a
measurable fraction of LAION-400M-class datasets for about sixty dollars, and
can frontrun the periodic snapshots of crowd-sourced sources such as Wikipedia
that RAG corpora are built from~\cite{carlini}. Standards now treat this as a
supply-chain risk: the OWASP Top 10 for LLM Applications lists \emph{Data and
Model Poisoning} (LLM04:2025) and mandates that deployments track data origins
and transformations, vet and sandbox unverified sources, and version their
data.\footnote{\href{https://github.com/OWASP/www-project-top-10-for-large-language-model-applications/blob/main/2_0_vulns/LLM04_DataModelPoisoning.md}{OWASP
Top 10 for LLM Applications, LLM04:2025 Data and Model Poisoning, 2025.}}

This raises a sharp operational question that existing work answers only in
part: once a source is found compromised, how quickly and how precisely can an
operator attribute the poisoned chunks to their source and purge every chunk
from that source? Query-time filters such as GMTP (a gradient-based masked-token
probability detector) remove poisoned documents at retrieval but keep no
provenance and offer no recovery~\cite{gmtp}. Traceback systems, RAGForensics
and RAGOrigin, attribute poisoned texts post hoc: triggered by a bad answer,
they re-retrieve candidate passages and prompt an LLM, one judging call per
candidate, to decide which are poisoned, an approach that is accurate but
reactive and costly per incident~\cite{ragforensics,ragorigin}. The closest
ingestion-time work, RAGShield, frames poisoning as a supply-chain problem and
attests content at admission, but attests at the \emph{document} level,
simulates its Ed25519 signature with a symmetric HMAC, and provides neither
signature-keyed traceback nor source revocation~\cite{ragshield}. A 2026 survey
confirms the gap: systematic support for provenance, rollback, and corpus
recovery remains under-developed~\cite{ragsurvey}.

We close this gap with \sysname{}, an ingestion gate that writes, per chunk, an
Ed25519-signed provenance record (source URI, principal, SHA-256 content hash,
detector verdicts, timestamp) natively into the vector store. Two side indices
turn the record into operational leverage: attributing a suspect chunk becomes a
single constant-time signature lookup instead of an iterative re-retrieval, and
revoking a compromised source becomes one command that enumerates and purges
exactly that source's chunks. On the real PoisonedRAG attack over Natural
Questions, with the \NsuspectsE{} top-ranked suspects surfaced by the released
dense-retrieval feedback, \sysname{} attributes poisoned chunks at recall and
precision comparable to the published RAGForensics LLM-judge baseline but in a
single lookup, about \textbf{\LatSpeedup$\times$} faster and at zero per-incident
model cost. The supply-chain framing and the ingestion gate are not themselves
novel; our contribution is the recovery layer that framing implies but that no
surveyed tool builds. Concretely:

\begin{itemize}
\item[\textbf{C1}] \textbf{Per-chunk signed provenance} that embeds the chunk's
  own content hash and admitting detector verdicts directly in the vector store,
  using \emph{real} Ed25519 public-key signatures rather than a symmetric
  stand-in (Section~\ref{sec:design}).
\item[\textbf{C2}] \textbf{Traceback as one constant-time signature lookup}, the
  structural alternative to iterative LLM-judge re-retrieval, evaluated against
  the published baseline on identical third-party inputs (Section~\ref{sec:eval}).
\item[\textbf{C3}] \textbf{One-command source revocation} that batch-purges
  every chunk of a compromised principal, the recovery layer absent from all
  surveyed tools, evaluated for mean-time-to-remediation at corpus scale
  (Section~\ref{sec:eval}).
\item[\textbf{C4}] A \textbf{reproducible open artifact} with a digest-pinned
  real corpus, third-party attack and baseline code, and an end-to-end
  evaluation harness (Sections~\ref{sec:eval} and~\ref{sec:limits}).
\end{itemize}

\section{Related Work}

We position \sysname{} against four lines of work and summarize the comparison
in Table~\ref{tab:related}.

\begin{table}[t]
\centering
\caption{Capabilities of prior work versus \sysname{}. Columns, left to right:
ingestion-time gate, per-chunk granularity, real public-key signatures,
signature-keyed traceback, one-command source revocation, and operational
evaluation on latency, recall, and mean-time-to-remediation.}
\label{tab:related}
\footnotesize
\setlength{\tabcolsep}{3.5pt}
\renewcommand{\arraystretch}{0.9}
\begin{tabular}{@{}lcccccc@{}}
\toprule
 & Ingest & Per- & Real & Sig.\ & Source & Oper.\\
Work & gate & chunk & sig. & traceback & revoke & eval \\
\midrule
RAGShield~\cite{ragshield}        & \cmark & \xmark & \xmark & \xmark & \xmark & \xmark \\
D-RAG (blockchain)                & \cmark & \xmark & \cmark & \xmark & \xmark & \xmark \\
RAGForensics~\cite{ragforensics}  & \xmark & \cmark & \xmark & \xmark & \xmark & \cmark \\
RAGOrigin~\cite{ragorigin}        & \xmark & \cmark & \xmark & \xmark & \xmark & \cmark \\
GMTP~\cite{gmtp}                  & \xmark & \cmark & \xmark & \xmark & \xmark & \xmark \\
RevPRAG~\cite{revprag}            & \xmark & \xmark & \xmark & \xmark & \xmark & \xmark \\
RAGPart/RAGMask~\cite{ragpartmask}& \xmark & \cmark & \xmark & \xmark & \xmark & \xmark \\
\textbf{\sysname{}}               & \cmark & \cmark & \cmark & \cmark & \cmark & \cmark \\
\bottomrule
\end{tabular}
\end{table}

\textbf{Detection and query-time filtering.} GMTP filters poisoned documents by
masking high-impact tokens and checking their probability under a masked
language model, eliminating over 90\% of poisoned content at
retrieval~\cite{gmtp}; RevPRAG detects poisoned responses at inference time from
LLM activations~\cite{revprag}; RAGPart and RAGMask are lightweight
retrieval-stage defenses that partition or flag retrieved
content~\cite{ragpartmask}. These contribute effective mitigation of the
\emph{effect} of an attack but keep no provenance, perform no attribution, and
offer no recovery; we treat them as a complementary query-time layer, not
competitors, and claim no novelty in detection. Indeed, on detector accuracy a
query-time filter like GMTP is the stronger tool; \sysname{} targets the
operational axis those tools do not address.

\textbf{Post-hoc traceback.} RAGForensics introduced traceback for RAG,
re-retrieving candidate texts and prompting an LLM to judge which are
poisoned~\cite{ragforensics}; RAGOrigin extends this to black-box responsibility
attribution with an adaptive scope and rank/relevance/influence
scoring~\cite{ragorigin}. Both are valuable forensic tools and report strong
attribution accuracy, but both are reactive, triggered per bad response, and
their iterative judging is costly at forensic time; neither attests provenance
at ingestion nor revokes a source. \sysname{} is complementary: it pays a small
one-time signing cost at ingestion so that attribution later becomes a
constant-time lookup, and adds the missing revocation step.

\textbf{Ingestion-time attestation.} RAGShield is the closest prior work, and we
credit it: it frames RAG poisoning as a supply-chain problem and attests content
at admission~\cite{ragshield}. Verified against its own text, however, it attests
at the document level, simulates its Ed25519 signature with a symmetric HMAC,
reports only attack-success and false-positive rates, and provides no
signature-keyed traceback or source revocation. D-RAG admits content with
provenance and human expert verification on a blockchain, but at document or
source granularity, with heavy machinery and no
one-command purge.\footnote{\href{https://aircconline.com/csit/papers/vol15/csit150417.pdf}{D-RAG:
A Privacy-Preserving Framework for Decentralized RAG using Blockchain, CSIT,
2025.}} The generic supply-chain analogues, namely Sigstore/Cosign signing and
the CycloneDX ML-BOM, sign at file granularity, carry no detector verdicts, and
are not instantiated for the fragments that RAG actually
retrieves.\footnote{\href{https://edu.chainguard.dev/open-source/sigstore/cosign/how-to-sign-an-sbom-with-cosign/}{How
to Sign an SBOM with Cosign, Chainguard Academy, 2024.}} \sysname{} is orthogonal
to all of these: per-chunk granularity, a real public-key signature, and the
recovery layer of traceback plus revocation that the analogy implies, as the
rightmost columns of Table~\ref{tab:related} make explicit.

\section{Design}\label{sec:design}

\sysname{} is an ingestion gate. For each chunk it computes a SHA-256 content
hash, runs a best-effort detector suite, assembles a canonical byte string over
the provenance tuple (chunk id, source URI, principal, content hash, detector
verdicts, timestamp, signer identity, granularity), signs it with the configured
backend, and writes the signed record into the datastore. The flow is:

\begin{center}\footnotesize
\fbox{\parbox{0.92\columnwidth}{\raggedright\setlength{\parindent}{0pt}%
\textbf{[1] Gate.} For each chunk: normalize and chunk; run detectors to obtain
verdicts; compute $h=\mathrm{SHA\text{-}256}(\text{chunk})$; sign
$\sigma=\mathrm{Ed25519}\bigl(k_{\mathrm{priv}},\ \text{URI}\,\|\,
\text{principal}\,\|\,h\,\|\,\text{verdicts}\,\|\,t\bigr)$.\par
\textbf{[2] Store.} Write the per-chunk record
$\langle\text{id},\text{URI},\text{principal},h,\text{verdicts},t,\sigma\rangle$
into the vector store; normal RAG queries are unchanged.\par
\textbf{[3] Traceback.} A suspect chunk is resolved by one constant-time
signature/hash lookup.\par
\textbf{[4] Revoke.} \texttt{revoke-source <principal>} batch-purges every chunk
of that principal.}}
\end{center}

\textbf{Why this is sound.} The datastore maintains two side indices. A
content-hash index maps a chunk hash to its chunk id, resolving a suspect chunk
to its signed record in constant time. A principal index maps a principal to its
set of chunk ids, so revoking a source enumerates exactly its chunks without
scanning the corpus. Signatures use real Ed25519, so provenance is
non-repudiable and verifiable by any holder of the public key. The signed
message binds the verdicts and the content hash, so a record cannot be replayed
for different content or different admitting verdicts without invalidating the
signature.

\textbf{Threat model and scope.} The attacker controls or compromises a web
source ingested by the corpus and injects poisoned texts, including the
split-view/frontrunning case where a source mutates after
collection~\cite{poisonedrag,carlini}. The protected asset is the integrity and
\emph{traceability} of the indexed corpus, and the secondary asset is
mean-time-to-remediation after a source is found compromised. The operator is
trusted and holds the private key; ingestion is where provenance is sealed.
\sysname{} does not claim to detect all poisoning at ingestion: detectors are
best-effort and may have false negatives, so the value is forensic and
operational, not perfect detection. Query-time detectors~\cite{gmtp,ragpartmask}
are a complementary layer.

\textbf{Reproducibility.} Behaviour is configured entirely by environment
variables (prefix \texttt{RAGTRAP\_}); nothing is hardcoded. Every run logs to
the console and to a timestamped file, recording the resolved configuration,
each input with its content hash, every step, and all outputs. A manifest pins
each third-party input by content digest, including the BEIR \texttt{nq} corpus,
the released attack-feedback file, and the PoisonedRAG passages by their
SHA-256, alongside model identities. The private key is generated per run and
never written to disk.

\section{Evaluation}\label{sec:eval}

\textbf{Setup and data.} The evaluation is non-circular by construction: the
attack, the retrieval, the ground-truth labels, and the baseline are all
third-party artifacts that \sysname{} did not author. The clean substrate is the
BEIR Natural Questions (\texttt{nq}) passage corpus, the exact
\FullPassages{}-passage Wikipedia corpus PoisonedRAG and RAGOrigin evaluate
on.\footnote{\href{https://huggingface.co/datasets/BeIR/nq}{BeIR/nq dataset
(Natural Questions passages), Hugging Face.}} The attack is PoisonedRAG's
released adversarial passages~\cite{poisonedrag}; the suspect set and its labels
come from the released RAGOrigin attack-feedback file (the top-100 contexts
surfaced by the \texttt{e5} dense retriever for each of \Nquestions{} target
questions, with third-party poisoned/clean
labels).\footnote{\href{https://github.com/zhangbl6618/RAG-Responsibility-Attribution}{RAG-Responsibility-Attribution
(RAGForensics/RAGOrigin code and released feedback), 2025.}} Both inputs are
pinned by SHA-256. We first validated the instrument (E0) on a labelled corpus
of \EZeroChunks{} chunks: all signed records verify, a tampered message is
rejected, hash-keyed traceback recovers the injected attribution, and revocation
purges exactly the targeted principal's \EZeroPurged{} chunks with no collateral.
We report E0 as a deterministic correctness guarantee, not a statistical
detection rate.

\textbf{E1: Traceback on the real attack (headline).} For each question we ingest
its retrieved contexts through the \sysname{} gate (real Ed25519 per chunk),
attributing poisoned contexts to a PoisonedRAG attacker principal and clean
contexts to benign sources. At forensic time the suspects are the top-\TopK{}
retrieved contexts per question (\NsuspectsE{} suspects: \NpoisonE{} poisoned,
\NcleanE{} clean), exactly what an attacked pipeline feeds the LLM. Two
attributors run on the \emph{identical} suspects: \sysname{}'s constant-time
content-hash lookup, and the published RAGForensics LLM-judge loop run verbatim
over a local \BLjudge{} model (one model call per context). As
Table~\ref{tab:e1} shows, \sysname{} attributes poisoned chunks at recall
\textbf{\RTrecall{}} and precision \textbf{\RTprecision{}} with false-positive
rate \RTfpr{}, comparable to the LLM-judge baseline's recall \BLrecall{}
(95\%~CI [\BLrecallLo{},\,\BLrecallHi{}]) and precision \BLprecision{}, but
\sysname{}'s lookup runs in \textbf{\RTlatPerSus~$\mu$s per suspect} versus
\BLperSus~s for the judge, about \textbf{\LatSpeedup$\times$} faster end to end
(\BLcalls{} model calls, $\sim$\BLcost~milli-USD at the published API rate,
versus zero model cost for \sysname{}). Precision is the property that is exact
by construction, since a content-hash match is exact; we frame it as a
correctness guarantee rather than a learned classifier's precision.

\begin{table}[t]
\centering
\caption{E1 traceback on the real PoisonedRAG attack over BEIR \texttt{nq}
(\Nquestions{} questions, top-\TopK{} suspects), \sysname{} versus the published
RAGForensics LLM-judge baseline on identical inputs. Proportions with 95\%
Wilson CIs; latency and cost are per incident.}
\label{tab:e1}
\footnotesize
\setlength{\tabcolsep}{5pt}
\renewcommand{\arraystretch}{0.9}
\begin{tabular}{@{}lcc@{}}
\toprule
Metric (\NsuspectsE{} suspects) & \sysname{} & RAGForensics judge \\
\midrule
Attribution recall   & \RTrecall{}$^{\dagger}$ & \BLrecall{} [\BLrecallLo{},\,\BLrecallHi{}] \\
Precision            & \RTprecision{}$^{\dagger}$ & \BLprecision{} [\BLprecisionLo{},\,\BLprecisionHi{}] \\
False-positive rate  & \RTfpr{}                & \BLfpr{} \\
Latency / suspect    & \textbf{\RTlatPerSus~$\mu$s} & \BLperSus~s \\
Work / suspect       & 1 lookup                & 1 model call \\
Per-incident cost    & \textbf{0}              & \BLcalls{} calls \\
\midrule
End-to-end speedup   & \multicolumn{2}{c}{\textbf{$\sim$\LatSpeedup$\times$}} \\
\bottomrule
\end{tabular}\\[2pt]
{\footnotesize $^{\dagger}$exact-hash recovery; a by-construction correctness
property for admitted chunks (see E1-drift).}
\end{table}

\textbf{E1-drift: honest false negatives under post-ingestion mutation.}
\sysname{} resolves a suspect only if its bytes were admitted through the gate, so
a poisoned chunk whose retrieved bytes drift after ingestion (a paraphrase, the
split-view case~\cite{carlini}) will not hash-match. We expose this directly:
paraphrasing a fraction of the poisoned suspects after ingestion drops recall to
\textbf{\RTrecallDriftThree{}} (95\%~CI [\RTrecallDriftThreeLo{},\,\RTrecallDriftThreeHi{}])
at 30\% drift and lower at 50\%, a real limitation we report rather than hide.
Precision stays exact (\RTprecision{}) because a hash match is exact; the missed
chunks are false negatives, not false positives. Recovering drifted variants
requires near-duplicate or semantic matching and is future work
(Section~\ref{sec:limits}).

\textbf{E2: Mean-time-to-remediation at corpus scale.} Treating one PoisonedRAG
attacker principal as compromised, we remove it two ways on the real corpus:
\sysname{}'s \texttt{revoke-source} batch purge versus a manual loop that scans
the whole corpus. \sysname{} revocation reads the principal's \RevChunks{}
chunk ids from the index and purges them in \textbf{\FullRevMttr~$\mu$s} on the
full \FullChunks{}-chunk corpus, while the manual scan takes \FullManMttr~ms, a
\textbf{\FullMttrRatio$\times$} advantage (Table~\ref{tab:scale}). The advantage
is structural: revocation cost scales with the revoked principal's chunk count
while the manual scan scales with the whole corpus, so the ratio grows with
corpus size, from \SmallMttrRatio$\times$ at \SmallPassages{} passages to
\FullMttrRatio$\times$ at the full corpus.

\begin{table}[t]
\centering
\caption{E2/E4 scaling on the real BEIR \texttt{nq} corpus. Per-chunk Ed25519
signing stays constant; the MTTR advantage of \texttt{revoke-source} (O(revoked))
over a full-corpus manual scan (O(corpus)) grows with corpus size.}
\label{tab:scale}
\footnotesize
\setlength{\tabcolsep}{5pt}
\renewcommand{\arraystretch}{0.9}
\begin{tabular}{@{}lcccc@{}}
\toprule
Passages & Chunks & Sign ($\mu$s) & Revoke & Manual / Revoke \\
\midrule
10k & 18k & 87 & 41~$\mu$s & \textbf{44$\times$} \\
100k & 172k & 86 & 52~$\mu$s & \textbf{627$\times$} \\
1M & 1.66M & 90 & 48~$\mu$s & \textbf{6{,}977$\times$} \\
2.68M & 4.36M & 69 & 46~$\mu$s & \textbf{16{,}931$\times$}%
\\
\bottomrule
\end{tabular}
\end{table}

\textbf{E3: Per-chunk versus per-document granularity.} On \EThreeDocs{} real
partially-poisoned documents (a real NQ passage into which real PoisonedRAG
passages are injected under one principal), the per-document scheme (sign the
parent, propagate trust) over-purges clean fragments: false-purge rate
\textbf{\PerDocFP{}} (95\%~CI [\PerDocFPLo{},\,\PerDocFPHi{}], \PerDocFalse{} of
\PerDocPurged{} removed chunks were clean), because it cannot distinguish clean
from poisoned fragments. \sysname{}'s per-chunk scheme localises the poisoned
chunks with false-purge rate \textbf{\PerChunkFP{}}.

\textbf{E4: Ingestion overhead and storage.} On the full corpus, signing each
chunk with real Ed25519 costs about \textbf{\EdLatency~$\mu$s}
($\sim$\EdThroughput{} chunks/s, single-threaded) and produces a fixed
\EdSig{}-byte signature in a record of about \EdRecord{} bytes. This is about
\textbf{\EdOverHmac$\times$} the cost of the symmetric HMAC stand-in
(\HmacLatency~$\mu$s, \HmacSig{}-byte tag); in exchange, public-key signing
yields non-repudiable provenance verifiable without the secret key, the
capability the simulated HMAC of RAGShield does not provide~\cite{ragshield}.
The overhead is small enough for practical ingestion.

\textbf{E5: Attack-success context.} To confirm the suspects are genuinely
dangerous, we feed the top-\TopK{} retrieved contexts to a local \AsrGen{}
generation model and check the answer: the attack steers it to the attacker's
target answer in \textbf{\AsrPoint\%} of the \AsrN{} questions (95\%~CI
[\AsrLo{},\,\AsrHi{}]). This is positioning context for the threat, not the
headline; it confirms that the chunks \sysname{} attributes do mislead the
pipeline it protects.

\section{Limitations and Conclusion}\label{sec:limits}

\textbf{Threats to validity.} External validity: E1's suspects and labels are
from one released attack-feedback set over NQ; the corpus-scale results (E2/E4)
use the full \FullPassages{}-passage corpus, and the structural arguments
(constant work for lookup and revocation) do not depend on scale. The baseline
judge is the published RAGForensics loop run over a local open model rather than
the hosted model the paper used; we report our measured numbers and price the
calls at the published API rate rather than claiming the original's scores.
Coverage: \sysname{}'s exact-hash attribution misses post-ingestion byte drift
(E1-drift); near-duplicate and semantic matching of suspects, adaptive and
multi-attacker regimes, and other corpora (HotpotQA, MS-MARCO, whose released
attack data we also obtained) are future work. Detection itself is out of scope
and best-effort, with query-time filters as the complementary
layer~\cite{gmtp,ragpartmask}. Key management: we assume a trusted operator
holding the private key, and key compromise is a residual risk.

\textbf{Conclusion.} Corpus poisoning is cheap and effective, standards now
demand provenance and quarantine, and the closest tools either filter at query
time, attribute post hoc with a costly iterative judge, or attest documents with
a simulated signature and no recovery step. The operational question we opened,
how fast and how precisely can an operator attribute and purge a compromised
source, has a concrete answer: by sealing per-chunk Ed25519 provenance at
ingestion, \sysname{} matches the published traceback baseline on attribution
but turns it into a constant-time lookup about \LatSpeedup$\times$ faster at zero
per-incident model cost, and turns recovery into one command whose advantage over
a full-corpus purge grows to \FullMttrRatio$\times$ at the full
\FullPassages{}-passage corpus, at an ingestion cost of \EdLatency~$\mu$s/chunk.
This is the recovery layer the supply-chain analogy implies but that no surveyed
tool had built. We release \sysname{} as a reproducible, double-blind
artifact.\footnote{\href{https://anonymous.4open.science}{Anonymized artifact
repository (link provided for double-blind review).}}

\bibliographystyle{sbc}
\bibliography{paper}

\end{document}
